\chapter{Funcionalidad y endpoints}
\label{ap:funcionalidad-endpoints}

La interfaz web de F.A.M.A. utiliza la API HTTP del backend para gestionar
sesiones, audios, entrenamientos y predicciones. Cada operación se
identifica por su método y su ruta. Cuando una tabla indica un prefijo, la
ruta completa se obtiene al concatenarlo con la ruta relativa. Las llaves
identifican parámetros que se sustituyen por el nombre o identificador del
recurso.

Las consultas utilizan \texttt{GET}. Las solicitudes que inician una
operación o registran información utilizan \texttt{POST}. Las rutas
\texttt{PUT}, \texttt{PATCH} y \texttt{DELETE} se emplean para transferencias,
asignaciones, modificaciones parciales o revocaciones según el recurso. La
respuesta de una operación en segundo plano se interpreta junto con el
estado del trabajo que devuelve la API.

\section{Estado del nodo e instalación}

La consulta \texttt{/health} comprueba que el proceso atiende solicitudes.
La consulta \texttt{/ready} informa si el nodo puede ofrecer el servicio
operativo según su configuración y esquema. El asistente se monta cuando
existe un directorio de configuración. Sus operaciones de detalle y
escritura requieren una autorización de instalación.

\begin{table}[!htbp]
    \centering
    \caption{Estado del nodo y configuración inicial.}
    \small
    \setlength{\tabcolsep}{4pt}
    \renewcommand{\arraystretch}{1.13}
    \begin{tabular}{@{}p{0.13\linewidth}p{0.36\linewidth}p{0.45\linewidth}@{}}
        \hline
        \textbf{Método} & \textbf{Ruta} & \textbf{Funcionalidad} \\
        \hline
        \texttt{GET} & \texttt{/\allowbreak{}} & Devuelve un mensaje de estado del backend operativo. \\
        \texttt{GET} & \texttt{/\allowbreak{}health} & Indica que el proceso HTTP responde. \\
        \texttt{GET} & \texttt{/\allowbreak{}ready} & Informa si el nodo está listo según el estado de configuración y esquema. \\
        \texttt{GET} & \texttt{/\allowbreak{}api/\allowbreak{}installation/\allowbreak{}status} & Consulta el estado general de la instalación. \\
        \texttt{GET} & \texttt{/\allowbreak{}api/\allowbreak{}installation/\allowbreak{}state} & Consulta el detalle del asistente con autorización de instalación. \\
        \texttt{POST} & \texttt{/\allowbreak{}api/\allowbreak{}installation/\allowbreak{}configuration} & Guarda la configuración del asistente autorizado. \\
        \texttt{POST} & \texttt{/\allowbreak{}api/\allowbreak{}installation/\allowbreak{}check} & Comprueba la configuración y las condiciones necesarias para continuar. \\
        \texttt{POST} & \texttt{/\allowbreak{}api/\allowbreak{}installation/\allowbreak{}complete} & Finaliza o reanuda el alta inicial y devuelve el estado de instalación. \\
        \hline
    \end{tabular}
\end{table}

\section{Acceso, cuentas y permisos}

La sesión del navegador se establece mediante una cookie \texttt{HttpOnly}.
Las operaciones de administración requieren una cuenta administradora.
Las rutas del catálogo de datasets requieren sesión y permiten al
investigador conocer el estado de su acceso o presentar una solicitud.
Las decisiones y asignaciones se administran mediante las rutas del
prefijo \texttt{/api/access}.

\begin{table}[!htbp]
    \centering
    \caption{Autenticación y gestión de la sesión.}
    \small
    \setlength{\tabcolsep}{4pt}
    \renewcommand{\arraystretch}{1.13}
    \begin{tabular}{@{}p{0.13\linewidth}p{0.36\linewidth}p{0.45\linewidth}@{}}
        \hline
        \multicolumn{3}{@{}l@{}}{\textbf{Prefijo} \texttt{/\allowbreak{}api/\allowbreak{}auth}} \\
        \hline
        \textbf{Método} & \textbf{Ruta relativa} & \textbf{Funcionalidad} \\
        \hline
        \texttt{POST} & \texttt{/\allowbreak{}login} & Recibe correo y contraseña. Establece la sesión web y devuelve el perfil. \\
        \texttt{GET} & \texttt{/\allowbreak{}me} & Consulta el perfil de la sesión vigente. \\
        \texttt{POST} & \texttt{/\allowbreak{}logout} & Cierra la sesión web y responde con HTTP 204. \\
        \texttt{POST} & \texttt{/\allowbreak{}change-password} & Cambia la contraseña de la cuenta autenticada e invalida su acceso de sesión. \\
        \texttt{POST} & \texttt{/\allowbreak{}reset-password} & Restablece la contraseña mediante una credencial temporal de recuperación. \\
        \texttt{POST} & \texttt{/\allowbreak{}register} & Crea una cuenta de investigador o administrador. Requiere rol administrador. \\
        \hline
    \end{tabular}
\end{table}

\begin{table}[!htbp]
    \centering
    \caption{Cuentas y seguridad.}
    \small
    \setlength{\tabcolsep}{4pt}
    \renewcommand{\arraystretch}{1.13}
    \begin{tabular}{@{}p{0.13\linewidth}p{0.36\linewidth}p{0.45\linewidth}@{}}
        \hline
        \multicolumn{3}{@{}l@{}}{\textbf{Prefijo} \texttt{/\allowbreak{}api}} \\
        \hline
        \textbf{Método} & \textbf{Ruta relativa} & \textbf{Funcionalidad} \\
        \hline
        \texttt{GET} & \texttt{/\allowbreak{}users} & Lista cuentas con paginación. \\
        \texttt{PATCH} & \texttt{/\allowbreak{}users/\allowbreak{}\{user\_\allowbreak{}id\}} & Actualiza el nombre, el rol o el estado de una cuenta. \\
        \texttt{POST} & \texttt{/\allowbreak{}users/\allowbreak{}\{user\_\allowbreak{}id\}/\allowbreak{}password-reset} & Emite un enlace temporal de recuperación de contraseña. \\
        \texttt{GET} & \texttt{/\allowbreak{}security/\allowbreak{}audit} & Consulta los eventos de seguridad mediante una lista paginada. \\
        \texttt{POST} & \texttt{/\allowbreak{}security/\allowbreak{}cleanup} & Retira registros expirados o fuera de retención en un lote acotado. \\
        \hline
    \end{tabular}
\end{table}

\begin{table}[!htbp]
    \centering
    \caption{Catálogo, solicitudes y asignaciones.}
    \small
    \setlength{\tabcolsep}{4pt}
    \renewcommand{\arraystretch}{1.13}
    \begin{tabular}{@{}p{0.13\linewidth}p{0.36\linewidth}p{0.45\linewidth}@{}}
        \hline
        \multicolumn{3}{@{}l@{}}{\textbf{Prefijo} \texttt{/\allowbreak{}api}} \\
        \hline
        \textbf{Método} & \textbf{Ruta relativa} & \textbf{Funcionalidad} \\
        \hline
        \texttt{GET} & \texttt{/\allowbreak{}datasets/\allowbreak{}catalogue} & Lista datasets registrados y el estado de acceso del usuario. \\
        \texttt{POST} & \texttt{/\allowbreak{}datasets/\allowbreak{}\{dataset\_\allowbreak{}id\}/\allowbreak{}requests} & Registra una solicitud de acceso para el investigador. \\
        \texttt{GET} & \texttt{/\allowbreak{}access/\allowbreak{}requests} & Lista solicitudes y permite filtrar su estado. \\
        \texttt{POST} & \texttt{/\allowbreak{}access/\allowbreak{}requests/\allowbreak{}\{user\_\allowbreak{}id\}/\allowbreak{}\{dataset\_\allowbreak{}id\}/\allowbreak{}decision} & Aprueba o rechaza una solicitud. La aprobación concede la asignación correspondiente. \\
        \texttt{GET} & \texttt{/\allowbreak{}access/\allowbreak{}users/\allowbreak{}\{user\_\allowbreak{}id\}/\allowbreak{}datasets} & Lista los datasets registrados y las asignaciones del usuario indicado. \\
        \texttt{GET} & \texttt{/\allowbreak{}access/\allowbreak{}users/\allowbreak{}\{user\_\allowbreak{}id\}/\allowbreak{}datasets/\allowbreak{}\{dataset\_\allowbreak{}id\}} & Consulta una asignación específica. \\
        \texttt{PUT} & \texttt{/\allowbreak{}access/\allowbreak{}users/\allowbreak{}\{user\_\allowbreak{}id\}/\allowbreak{}datasets/\allowbreak{}\{dataset\_\allowbreak{}id\}} & Concede acceso al dataset indicado. \\
        \texttt{DELETE} & \texttt{/\allowbreak{}access/\allowbreak{}users/\allowbreak{}\{user\_\allowbreak{}id\}/\allowbreak{}datasets/\allowbreak{}\{dataset\_\allowbreak{}id\}} & Revoca la asignación y responde con HTTP 204. \\
        \hline
    \end{tabular}
\end{table}

\section{Gestión y preparación de conjuntos de datos}

La API de ingesta permite consultar el inventario, cargar audios
etiquetados, descargar desde GCS y preprocesar una selección local. El
investigador trabaja sobre datasets asignados y el administrador puede
gestionar el inventario general. El diagnóstico de conexión del asistente
y el estado de ingesta se relacionan con \texttt{RF\_01}. La sincronización
y el preprocesamiento corresponden a \texttt{RF\_02} y \texttt{RF\_03}.

Las rutas de particiones permiten definir y recuperar una selección
versionada de grabaciones para entrenamiento, validación y prueba. Guardar
una versión fija la preparación que podrá seleccionarse al iniciar un
trabajo.

\begin{table}[!htbp]
    \centering
    \caption{Ingesta y preparación de audios.}
    \small
    \setlength{\tabcolsep}{4pt}
    \renewcommand{\arraystretch}{1.13}
    \begin{tabular}{@{}p{0.13\linewidth}p{0.36\linewidth}p{0.45\linewidth}@{}}
        \hline
        \multicolumn{3}{@{}l@{}}{\textbf{Prefijo} \texttt{/\allowbreak{}api/\allowbreak{}ingestion}} \\
        \hline
        \textbf{Método} & \textbf{Ruta relativa} & \textbf{Funcionalidad} \\
        \hline
        \texttt{GET} & \texttt{/\allowbreak{}status} & Consulta el estado de GCS y el inventario general. Requiere rol administrador. \\
        \texttt{GET} & \texttt{/\allowbreak{}datasets} & Lista los datasets visibles y los límites de carga disponibles. \\
        \texttt{GET} & \texttt{/\allowbreak{}datasets/\allowbreak{}\{dataset\_\allowbreak{}name\}/\allowbreak{}files} & Consulta el inventario y los metadatos disponibles del dataset. \\
        \texttt{POST} & \texttt{/\allowbreak{}upload} & Recibe archivos y su destino mediante un formulario multipart. Carga los audios aceptados y registra sus metadatos. \\
        \texttt{POST} & \texttt{/\allowbreak{}sync} & Sincroniza audios desde GCS al nodo. Puede solicitar preprocesamiento y devuelve el resultado por dataset. \\
        \texttt{POST} & \texttt{/\allowbreak{}preprocess} & Genera tensores a partir de los audios locales seleccionados y devuelve los totales procesados y fallidos. \\
        \hline
    \end{tabular}
\end{table}

\begin{table}[!htbp]
    \centering
    \caption{Particiones versionadas.}
    \small
    \setlength{\tabcolsep}{4pt}
    \renewcommand{\arraystretch}{1.13}
    \begin{tabular}{@{}p{0.13\linewidth}p{0.36\linewidth}p{0.45\linewidth}@{}}
        \hline
        \multicolumn{3}{@{}l@{}}{\textbf{Prefijo} \texttt{/\allowbreak{}api/\allowbreak{}training/\allowbreak{}datasets/\allowbreak{}\{dataset\_\allowbreak{}name\}/\allowbreak{}partitions}} \\
        \hline
        \textbf{Método} & \textbf{Ruta relativa} & \textbf{Funcionalidad} \\
        \hline
        \texttt{GET} & \emph{Prefijo exacto} & Lista las versiones disponibles para el dataset. \\
        \texttt{POST} & \texttt{/\allowbreak{}preview} & Calcula una propuesta de reparto con sus cantidades y advertencias. \\
        \texttt{POST} & \emph{Prefijo exacto} & Guarda una versión a partir de una propuesta identificada. \\
        \texttt{GET} & \texttt{/\allowbreak{}\{version\}} & Consulta el manifiesto de una versión de particiones. \\
        \hline
    \end{tabular}
\end{table}

\section{Entrenamiento y gestión de modelos}

La operación de inicio recibe el dataset, los parámetros de entrenamiento
y la identidad de la preparación cuando corresponde. Su respuesta indica
el resultado de la admisión. La interfaz consulta el progreso para conocer
la evolución y el estado final del trabajo, en el marco de
\texttt{RF\_04}. El investigador puede consultar o detener su propio
trabajo mientras conserva acceso al dataset. El administrador puede
supervisar el trabajo del nodo.

La activación se realiza mediante una operación administrativa explícita.
El historial, las evaluaciones y las copias del checkpoint permiten
revisar y recuperar un modelo. Las operaciones de copia requieren acceso
al modelo y una identidad de artefacto verificable.

\begin{table}[!htbp]
    \centering
    \caption{Entrenamiento y seguimiento.}
    \small
    \setlength{\tabcolsep}{4pt}
    \renewcommand{\arraystretch}{1.13}
    \begin{tabular}{@{}p{0.13\linewidth}p{0.36\linewidth}p{0.45\linewidth}@{}}
        \hline
        \multicolumn{3}{@{}l@{}}{\textbf{Prefijo} \texttt{/\allowbreak{}api/\allowbreak{}training}} \\
        \hline
        \textbf{Método} & \textbf{Ruta relativa} & \textbf{Funcionalidad} \\
        \hline
        \texttt{GET} & \texttt{/\allowbreak{}hardware} & Consulta la información disponible de CPU, GPU y memoria. \\
        \texttt{GET} & \texttt{/\allowbreak{}datasets} & Lista los datasets accesibles que el servicio reconoce como fuentes de entrenamiento. \\
        \texttt{POST} & \texttt{/\allowbreak{}start} & Valida la fuente y los parámetros e inicia el trabajo en segundo plano. \\
        \texttt{GET} & \texttt{/\allowbreak{}progress} & Devuelve estado, avance por época, métricas y registros del trabajo. \\
        \texttt{POST} & \texttt{/\allowbreak{}stop} & Solicita la detención cooperativa del entrenamiento activo. \\
        \texttt{GET} & \texttt{/\allowbreak{}history} & Consulta el historial de modelos entrenados visible para la sesión. \\
        \texttt{POST} & \texttt{/\allowbreak{}models/\allowbreak{}\{model\_\allowbreak{}id\}/\allowbreak{}activate} & Solicita la designación del modelo predeterminado. Requiere rol administrador. \\
        \hline
    \end{tabular}
\end{table}

\begin{table}[!htbp]
    \centering
    \caption{Copias de checkpoints.}
    \small
    \setlength{\tabcolsep}{4pt}
    \renewcommand{\arraystretch}{1.13}
    \begin{tabular}{@{}p{0.13\linewidth}p{0.36\linewidth}p{0.45\linewidth}@{}}
        \hline
        \multicolumn{3}{@{}l@{}}{\textbf{Prefijo} \texttt{/\allowbreak{}api/\allowbreak{}training/\allowbreak{}models/\allowbreak{}\{model\_\allowbreak{}id\}/\allowbreak{}checkpoint}} \\
        \hline
        \textbf{Método} & \textbf{Ruta relativa} & \textbf{Funcionalidad} \\
        \hline
        \texttt{GET} & \emph{Prefijo exacto} & Consulta el estado de las copias local y remota. \\
        \texttt{POST} & \texttt{/\allowbreak{}backup} & Respalda el checkpoint local en GCS y verifica su identidad. \\
        \texttt{POST} & \texttt{/\allowbreak{}restore} & Recupera y verifica una copia local a partir del respaldo de GCS. \\
        \hline
    \end{tabular}
\end{table}

\section{Predicción y retroalimentación}

La predicción recibe un audio WAV mediante un formulario multipart. El
identificador opcional del modelo se envía como parámetro de consulta.
Cuando se completa la operación, la respuesta incluye el identificador
del registro en PostgreSQL, la clase, la puntuación de confianza y la
identidad del modelo utilizado. Esta operación se relaciona con
\texttt{RF\_05}.

Para \texttt{RF\_06}, la interfaz utiliza el identificador de predicción
para registrar la confirmación o corrección. El autor de la
retroalimentación se obtiene de la sesión. La incorporación local y la
sincronización posterior con GCS tienen estados separados y se consultan
mediante las rutas de revisión. El acceso a las predicciones y a sus
aportes se limita a los recursos autorizados.

\begin{table}[!htbp]
    \centering
    \caption{Predicción, catálogo y estadísticas.}
    \small
    \setlength{\tabcolsep}{4pt}
    \renewcommand{\arraystretch}{1.13}
    \begin{tabular}{@{}p{0.13\linewidth}p{0.36\linewidth}p{0.45\linewidth}@{}}
        \hline
        \multicolumn{3}{@{}l@{}}{\textbf{Prefijo} \texttt{/\allowbreak{}api}} \\
        \hline
        \textbf{Método} & \textbf{Ruta relativa} & \textbf{Funcionalidad} \\
        \hline
        \texttt{GET} & \texttt{/\allowbreak{}models} & Devuelve modelos visibles, clases, disponibilidad de pesos y evaluaciones registradas cuando están disponibles. \\
        \texttt{GET} & \texttt{/\allowbreak{}model-info} & Consulta el estado observado del servicio bioacústico histórico. \\
        \texttt{GET} & \texttt{/\allowbreak{}dashboard/\allowbreak{}stats} & Consulta las estadísticas operativas visibles según el acceso de la sesión. \\
        \texttt{POST} & \texttt{/\allowbreak{}predict} & Recibe un WAV y un identificador opcional de modelo. Devuelve la clase, la confianza y el identificador de predicción. \\
        \hline
    \end{tabular}
\end{table}

\begin{table}[!htbp]
    \centering
    \caption{Retroalimentación y curación.}
    \small
    \setlength{\tabcolsep}{4pt}
    \renewcommand{\arraystretch}{1.13}
    \begin{tabular}{@{}p{0.13\linewidth}p{0.36\linewidth}p{0.45\linewidth}@{}}
        \hline
        \multicolumn{3}{@{}l@{}}{\textbf{Prefijo} \texttt{/\allowbreak{}api/\allowbreak{}feedback}} \\
        \hline
        \textbf{Método} & \textbf{Ruta relativa} & \textbf{Funcionalidad} \\
        \hline
        \texttt{POST} & \emph{Prefijo exacto} & Registra la confirmación o corrección de una predicción. Devuelve el identificador de la retroalimentación. \\
        \texttt{GET} & \texttt{/\allowbreak{}pending} & Consulta los aportes pendientes de revisión. \\
        \texttt{POST} & \texttt{/\allowbreak{}\{id\_\allowbreak{}retroalimentacion\}/\allowbreak{}approve} & Incorpora el aporte al dataset local y registra su estado de sincronización. \\
        \texttt{POST} & \texttt{/\allowbreak{}\{id\_\allowbreak{}retroalimentacion\}/\allowbreak{}reject} & Registra el descarte del aporte de curación. \\
        \texttt{GET} & \texttt{/\allowbreak{}sync} & Consulta el estado de sincronización de los aportes incorporados. \\
        \texttt{GET} & \texttt{/\allowbreak{}stats} & Devuelve totales y distribución de validaciones y correcciones. \\
        \hline
    \end{tabular}
\end{table}

\section{Importación y exportación ZIP}

El módulo ZIP se habilita en entornos POSIX con los procesos auxiliares
configurados. En Windows nativo, la consulta de capacidad informa que el
módulo está indisponible. El acceso a una importación o a un dataset
importado se restringe al titular del recurso. La descarga de una exportación utiliza su propia
autorización temporal.

La importación puede recibir el ZIP completo o mediante fragmentos.
Las rutas de previsualización, confirmación y sincronización admiten tareas
en segundo plano. Una respuesta HTTP 202 indica que el trabajo fue
aceptado. El cliente consulta la ruta de estado de la ejecución para
conocer el resultado. La exportación también se solicita para una revisión
identificada y se consulta hasta que el ZIP esté disponible.

\begin{table}[!htbp]
    \centering
    \caption{Importación ZIP y recepción de archivos.}
    \small
    \setlength{\tabcolsep}{4pt}
    \renewcommand{\arraystretch}{1.13}
    \begin{tabular}{@{}p{0.13\linewidth}p{0.36\linewidth}p{0.45\linewidth}@{}}
        \hline
        \multicolumn{3}{@{}l@{}}{\textbf{Prefijo} \texttt{/\allowbreak{}api/\allowbreak{}ingestion/\allowbreak{}imports}} \\
        \hline
        \textbf{Método} & \textbf{Ruta relativa} & \textbf{Funcionalidad} \\
        \hline
        \texttt{GET} & \texttt{/\allowbreak{}access} & Consulta la disponibilidad del módulo y el estado de acceso. \\
        \texttt{GET} & \texttt{/\allowbreak{}limits} & Consulta los límites de la importación. \\
        \texttt{GET} & \texttt{/\allowbreak{}example.zip} & Descarga un archivo de ejemplo. \\
        \texttt{POST} & \emph{Prefijo exacto} & Crea el registro de una importación. \\
        \texttt{GET} & \emph{Prefijo exacto} & Lista las importaciones de la identidad autorizada. \\
        \texttt{GET} & \texttt{/\allowbreak{}\{job\_\allowbreak{}id\}} & Consulta el estado del registro de importación. \\
        \texttt{PUT} & \texttt{/\allowbreak{}\{job\_\allowbreak{}id\}/\allowbreak{}archive} & Recibe el ZIP completo como contenido binario de tipo application/zip. \\
        \texttt{POST} & \texttt{/\allowbreak{}\{job\_\allowbreak{}id\}/\allowbreak{}receipt} & Reserva la recepción de un archivo para la importación. \\
        \texttt{GET} & \texttt{/\allowbreak{}\{job\_\allowbreak{}id\}/\allowbreak{}receipt} & Consulta el estado de la recepción. \\
        \texttt{PUT} & \texttt{/\allowbreak{}\{job\_\allowbreak{}id\}/\allowbreak{}receipt/\allowbreak{}chunks/\allowbreak{}\{index\}} & Envía un fragmento con su desplazamiento y huella. \\
        \texttt{GET} & \texttt{/\allowbreak{}\{job\_\allowbreak{}id\}/\allowbreak{}receipt/\allowbreak{}chunks} & Consulta los fragmentos confirmados. \\
        \texttt{POST} & \texttt{/\allowbreak{}\{job\_\allowbreak{}id\}/\allowbreak{}receipt/\allowbreak{}finalize} & Finaliza la recepción por fragmentos. \\
        \texttt{POST} & \texttt{/\allowbreak{}\{job\_\allowbreak{}id\}/\allowbreak{}receipt/\allowbreak{}cancel} & Solicita cancelar la recepción asociada. \\
        \hline
    \end{tabular}
\end{table}

\begin{table}[!htbp]
    \centering
    \caption{Revisión y ejecución de importaciones ZIP.}
    \small
    \setlength{\tabcolsep}{4pt}
    \renewcommand{\arraystretch}{1.13}
    \begin{tabular}{@{}p{0.13\linewidth}p{0.36\linewidth}p{0.45\linewidth}@{}}
        \hline
        \multicolumn{3}{@{}l@{}}{\textbf{Prefijo} \texttt{/\allowbreak{}api/\allowbreak{}ingestion/\allowbreak{}imports/\allowbreak{}\{job\_\allowbreak{}id\}}} \\
        \hline
        \textbf{Método} & \textbf{Ruta relativa} & \textbf{Funcionalidad} \\
        \hline
        \texttt{POST} & \texttt{/\allowbreak{}preview} & Admite la validación y previsualización del archivo. \\
        \texttt{GET} & \texttt{/\allowbreak{}files} & Consulta los archivos detectados durante la revisión. \\
        \texttt{GET} & \texttt{/\allowbreak{}classes} & Consulta las clases detectadas durante la revisión. \\
        \texttt{GET} & \texttt{/\allowbreak{}issues} & Consulta los problemas identificados durante la revisión. \\
        \texttt{POST} & \texttt{/\allowbreak{}confirm} & Admite la publicación de una importación a partir de su previsualización. \\
        \texttt{POST} & \texttt{/\allowbreak{}cancel} & Solicita cancelar la importación. \\
        \texttt{POST} & \texttt{/\allowbreak{}sync} & Admite una tarea de sincronización de la importación con GCS. \\
        \texttt{GET} & \texttt{/\allowbreak{}jobs/\allowbreak{}\{execution\_\allowbreak{}id\}} & Consulta el estado de una ejecución de la importación. \\
        \texttt{POST} & \texttt{/\allowbreak{}jobs/\allowbreak{}\{execution\_\allowbreak{}id\}/\allowbreak{}cancel} & Solicita cancelar una ejecución. \\
        \texttt{POST} & \texttt{/\allowbreak{}jobs/\allowbreak{}\{execution\_\allowbreak{}id\}/\allowbreak{}retry} & Admite el reintento de una ejecución fallida o cancelada con una nueva clave de solicitud. \\
        \hline
    \end{tabular}
\end{table}

\begin{table}[!htbp]
    \centering
    \caption{Datasets importados y exportación ZIP.}
    \small
    \setlength{\tabcolsep}{4pt}
    \renewcommand{\arraystretch}{1.13}
    \begin{tabular}{@{}p{0.13\linewidth}p{0.36\linewidth}p{0.45\linewidth}@{}}
        \hline
        \multicolumn{3}{@{}l@{}}{\textbf{Prefijo} \texttt{/\allowbreak{}api/\allowbreak{}ingestion/\allowbreak{}imports}} \\
        \hline
        \textbf{Método} & \textbf{Ruta relativa} & \textbf{Funcionalidad} \\
        \hline
        \texttt{GET} & \texttt{/\allowbreak{}datasets} & Lista los datasets importados visibles para el titular. \\
        \texttt{GET} & \texttt{/\allowbreak{}datasets/\allowbreak{}\{dataset\_\allowbreak{}id\}/\allowbreak{}files} & Lista los archivos de un dataset importado. \\
        \texttt{GET} & \texttt{/\allowbreak{}datasets/\allowbreak{}\{dataset\_\allowbreak{}id\}/\allowbreak{}files/\allowbreak{}\{sha256\}} & Descarga el WAV identificado por su huella. \\
        \texttt{GET} & \texttt{/\allowbreak{}datasets/\allowbreak{}\{dataset\_\allowbreak{}id\}/\allowbreak{}revisions/\allowbreak{}\{revision\}} & Consulta una revisión del dataset importado. \\
        \texttt{GET} & \texttt{/\allowbreak{}datasets/\allowbreak{}\{dataset\_\allowbreak{}id\}/\allowbreak{}revisions/\allowbreak{}\{revision\}/\allowbreak{}manifest} & Entrega el manifiesto de la revisión indicada. \\
        \texttt{POST} & \texttt{/\allowbreak{}datasets/\allowbreak{}\{dataset\_\allowbreak{}id\}/\allowbreak{}revisions/\allowbreak{}\{revision\}/\allowbreak{}exports} & Solicita una exportación ZIP para la revisión indicada. \\
        \texttt{GET} & \texttt{/\allowbreak{}datasets/\allowbreak{}\{dataset\_\allowbreak{}id\}/\allowbreak{}exports} & Consulta el historial de exportaciones del dataset. \\
        \texttt{GET} & \texttt{/\allowbreak{}exports/\allowbreak{}\{export\_\allowbreak{}id\}} & Consulta el estado de una exportación. \\
        \texttt{POST} & \texttt{/\allowbreak{}exports/\allowbreak{}\{export\_\allowbreak{}id\}/\allowbreak{}sessions} & Emite una sesión temporal para descargar el ZIP disponible. \\
        \texttt{GET/HEAD} & \texttt{/\allowbreak{}exports/\allowbreak{}\{export\_\allowbreak{}id\}/\allowbreak{}download} & Descarga el ZIP o consulta sus cabeceras con autorización de descarga. \\
        \hline
    \end{tabular}
\end{table}

\section{Interpretación de respuestas}

Los códigos HTTP y el estado de la operación se leen en conjunto.
HTTP 201 indica creación o registro en las operaciones que lo declaran.
HTTP 204 indica una respuesta exitosa sin cuerpo. HTTP 401 corresponde a
credenciales o sesión inválidas y HTTP 403 a una operación denegada.
HTTP 404 se utiliza para recursos inexistentes o fuera del acceso
permitido. HTTP 409 informa un conflicto con el estado del recurso y
HTTP 422 indica una solicitud que no cumple la validación. Las cargas
pueden responder con HTTP 413 cuando exceden el límite y con HTTP 415
cuando el tipo de contenido es incompatible con la ruta. HTTP 429
señala que se alcanzó un límite de solicitudes. HTTP 503 indica
indisponibilidad según el servicio, el modelo o la capacidad solicitada.

En los trabajos de entrenamiento y las transferencias, la interfaz
utiliza los campos de estado de la respuesta para distinguir el inicio,
el avance, la finalización y los errores parciales.

% Rutas contrastadas con las declaraciones del backend.
% backend/app/entrypoint.py
% backend/app/installation.py
% backend/app/main.py
% backend/app/routes/auth.py
% backend/app/routes/users.py
% backend/app/routes/security.py
% backend/app/routes/access.py
% backend/app/routes/partitions.py
% backend/app/routes/checkpoints.py
% backend/app/routes/import_capability.py
% backend/app/routes/dataset_imports.py
% backend/app/routes/dataset_exports.py
